\documentclass[a4paper,12pt]{article}

% 页面设置
\usepackage[top=2.54cm, bottom=2.54cm, left=1.91cm, right=1.91cm]{geometry}

% 中文支持
\usepackage[UTF8]{ctex}
\usepackage{fontspec}

% 数学公式支持
\usepackage{amsmath}
\usepackage{amssymb}
\usepackage{amsfonts}
\usepackage{amsthm}

\usepackage{enumitem}

% 定理环境定义
\theoremstyle{definition}
\newtheorem*{pf}{证}
\newtheorem*{sol}{解}

% 图形支持
\usepackage{graphicx}
\usepackage{tikz}

% 超链接支持
\usepackage{hyperref}
\hypersetup{
    colorlinks=true,
    linkcolor=blue,
    filecolor=magenta,
    urlcolor=cyan,
}

% 行距设置
\linespread{1.25}

% 标题信息
\title{Mathematical Analysis II\\Week 8 Homework (April 23)}
\author{袁子强\hspace{1cm} 2025533009}
% \date{}

\begin{document}

\maketitle

\section*{Section 10.2}

1. 计算下列积分.

(1) $\displaystyle\int_{0}^{R}\,\mathrm{d}x \int_{0}^{\sqrt{R^2 - x^2}} \ln(1 + x^2 + y^2)\,\mathrm{d}y$;
\begin{sol}
    令$\begin{cases}
            x=r\cos\theta \\
            y=r\sin\theta
        \end{cases}$, 所以$\mathrm{d}x\mathrm{d}y=r\,\mathrm{d}r\mathrm{d}\theta$

    所以$\int_{0}^{R}\,\mathrm{d}x \int_{0}^{\sqrt{R^2 - x^2}} \ln(1 + x^2 + y^2)\,\mathrm{d}y=\int_0^{\frac{\pi}{2}}\,\mathrm{d}\theta\int_0^Rr\ln(1+r^2)\,\mathrm{d}r$

    其中
    \begin{align*}
        \int_0^Rr\ln(1+r^2)\,\mathrm{d}r
         & =\int_0^R\frac{1}{2}\ln(1+r^2)\,\mathrm{d}(1+r^2)            \\
         & \overset{u=1+r^2}{=}\frac{1}{2}\left.(u\ln u-u)\right|_{0}^R \\
         & =\frac{1}{2}\left.((1+r^2)\ln(1+r^2)-1-r^2)\right|_{0}^R     \\
         & =\frac{1}{2}((1+R^2)\ln(1+R^2)-R^2)
    \end{align*}

    所以
    \begin{align*}
        \int_{0}^{R}\,\mathrm{d}x \int_{0}^{\sqrt{R^2 - x^2}} \ln(1 + x^2 + y^2)\,\mathrm{d}y
         & =\int_0^{\frac{\pi}{2}}\frac{1}{2}((1+R^2)\ln(1+R^2)-R^2)\,\mathrm{d} \\
         & =\frac{\pi}{4}((1+R^2)\ln(1+R^2)-R^2)
    \end{align*}
\end{sol}
(3) $\displaystyle\int_{0}^{\pi} \int_{0}^{\pi} \cos(x + y)\,\mathrm{d}x\mathrm{d}y$;
\begin{sol}
    \begin{align*}
        \text{原式}
         & =\int_{0}^{\pi}\left.\sin(x+y)\right|_0^{\pi}\,\mathrm{d}y \\
         & =-2\int_{0}^{\pi}\sin y\,\mathrm{d}y                       \\
         & =2\left.(\cos y)\right|_0^{\pi}                            \\
         & =-4
    \end{align*}
\end{sol}
(5) $\displaystyle\int_{0}^{\frac{R}{\sqrt{1+R^2}}}\,\mathrm{d}x \int_{0}^{Rx} \left(1 + \frac{y^2}{x^2}\right)\,\mathrm{d}y + \int_{\frac{R}{\sqrt{1+R^2}}}^{R}\,\mathrm{d}x \int_{0}^{\sqrt{R^2 - x^2}} \left(1 + \frac{y^2}{x^2}\right)\,\mathrm{d}y$;
\begin{sol}
    即求$y=Rx$, $x^2+y^2=R^2$, $y=0$围成的扇形面积

    令$\begin{cases}
            x=r\cos\theta \\
            y=r\sin\theta
        \end{cases}$

    $\theta$范围：$0\leq\theta\leq\arctan R$

    所以$\begin{cases}
            1+\frac{y^2}{x^2}=1+\cot^2\theta=\sec^2\theta \\
            \mathrm{d}x\mathrm{d}y=r\,\mathrm{d}r\mathrm{d}\theta
        \end{cases}$

    所以
    \begin{align*}
        \text{原式}
         & =\int_0^{\arctan R}\,\mathrm{d}\theta\int_0^Rr\sec^2\theta\,\mathrm{d}r         \\
         & =\int_0^{\arctan R}\left.\frac{r^2}{2}\sec^2\theta\right|_0^R\,\mathrm{d}\theta \\
         & =\frac{R^2}{2}\int_0^{\arctan R}\sec^2\theta\,\mathrm{d}\theta                  \\
         & =\frac{R^2}{2}\left.(\tan\theta)\right|_0^{\arctan R}                           \\
         & =\frac{R^3}{2}
    \end{align*}
\end{sol}


2. 计算下面二重积分.

(1) $\displaystyle\iint_{D} \sqrt{x^2 + y^2}\,\mathrm{d}x\mathrm{d}y$, $D$: $x^2 + y^2 \leq x + y$;
\begin{sol}
    令$\begin{cases}
            x=r\cos\theta \\
            y=r\sin\theta
        \end{cases}$，所以$\mathrm{d}x\mathrm{d}y=r\,\mathrm{d}r\mathrm{d}\theta$, $0\leq r\leq x+y=(\cos\theta+\sin\theta)=\sqrt{2}\sin(\theta+\frac{\pi}{4})$

    所以$0\leq\theta+\frac{\pi}{4}\leq\pi$，所以$-\frac{\pi}{4}\leq\theta\leq\frac{3\pi}{4}$

    所以

    \begin{align*}
        \text{原式}
         & =\iint_{D}r^2\,\mathrm{d}r\mathrm{d}\theta                                                                           \\
         & =\int_{-\frac{\pi}{4}}^{\frac{3\pi}{4}}\,\mathrm{d}\theta\int_0^{\sqrt{2}\sin(\theta+\frac{\pi}{4})}r^2\,\mathrm{d}r \\
         & =\frac{2\sqrt{2}}{3}\int_{-\frac{\pi}{4}}^{\frac{3\pi}{4}}\sin^3(\theta+\frac{\pi}{4})\,\mathrm{d}\theta
    \end{align*}

    令$\alpha=\theta+\frac{\pi}{4}$，所以

    \begin{align*}
        \text{原式}
         & =\frac{2\sqrt{2}}{3}\int_0^{\pi}\sin^3\alpha\,\mathrm{d}\alpha           \\
         & =\frac{4\sqrt{2}}{3}\int_0^{\frac{\pi}{2}}\sin^3\alpha\,\mathrm{d}\alpha \\
         & =\frac{8\sqrt{2}}{9}
    \end{align*}
\end{sol}
(3) $\displaystyle\iint_{D} (x^2 + y^2)\,\mathrm{d}x\mathrm{d}y$, $D$: 由 $xy=1$, $xy=2$, $y=x$, $y=2x$ 围成的第一象限部分;
\begin{sol}
    令$\begin{cases}
            u=xy \\
            v=\frac{y}{x}
        \end{cases}$

    所以$\begin{cases}
            1\leq u\leq 2 \\
            1\leq v\leq 2
        \end{cases}$, $\begin{cases}
            x^2=\frac{u}{v} \\
            y^2=uv
        \end{cases}$, $\frac{\partial(u,v)}{\partial(x,y)}=\begin{vmatrix}
            y              & x           \\
            -\frac{y}{x^2} & \frac{1}{x}
        \end{vmatrix}=\frac{y}{x}+\frac{y}{x}=\frac{2y}{x}=2v$

    所以$\mathrm{d}x\mathrm{d}y=\frac{1}{2v}\,\mathrm{d}u\mathrm{d}v$

    所以

    \begin{align*}
        \text{原式}
         & =\int_1^2\,\mathrm{d}v\int_1^2\frac{u}{2v^2}(1+v^2)\,\mathrm{d}u \\
         & =\frac{3}{2}\int_1^2\frac{1}{2v^2}+\frac{1}{2}\,\mathrm{d}v      \\
         & =\frac{3}{4}\left.(v-\frac{1}{v})\right|_1^2                     \\
         & =\frac{9}{8}
    \end{align*}
\end{sol}
(5) $\displaystyle\iint_{D} xy\,\mathrm{d}x\mathrm{d}y$, $D$: $xy=a$, $xy=b$, $y^2=cx$, $y^2=dx$ 围成的第一象限部分 ($0 < a < b$, $0 < c < d$);
\begin{sol}
    令$\begin{cases}
            u=xy \\
            v=\frac{y^2}{x}
        \end{cases}$

    所以$\begin{cases}
            a\leq u\leq b \\
            c\leq v\leq d
        \end{cases}$, $\frac{\partial(u,v)}{\partial(x,y)}=\begin{vmatrix}
            y                & x            \\
            -\frac{y^2}{x^2} & \frac{2y}{x}
        \end{vmatrix}=\frac{2y^2}{x}+\frac{y^2}{x}=\frac{3y^2}{x}=3v$

    所以$\mathrm{d}x\mathrm{d}y=\frac{1}{3v}\,\mathrm{d}u\mathrm{d}v$

    所以

    \begin{align*}
        \text{原式}
         & =\int_c^d\,\mathrm{d}v\int_a^b\frac{u}{3v}\,\mathrm{d}u \\
         & =\frac{b^2-a^2}{6}\int_c^d\frac{1}{v}\,\mathrm{d}v      \\
         & =\frac{b^2-a^2}{6}\ln\frac{d}{c}
    \end{align*}
\end{sol}
(7) $\displaystyle\iint_{D} \frac{x^2 - y^2}{\sqrt{x + y + 3}}\,\mathrm{d}x\mathrm{d}y$, $D$: $|x| + |y| \leq 1$;
\begin{sol}
    令$\begin{cases}
            u=x+y \\
            v=x-y
        \end{cases}$

    所以$\begin{cases}
            -1\leq u\leq 1 \\
            -1\leq v\leq 1
        \end{cases}$, $\begin{cases}
            x^2-y^2=uv \\
            \sqrt{x+y+3}=\sqrt{u+3}
        \end{cases}$, $\frac{\partial(u,v)}{\partial(x,y)}=\begin{vmatrix}
            1 & 1  \\
            1 & -1
        \end{vmatrix}=-2$

    所以$\mathrm{d}x\mathrm{d}y=\frac{1}{2}\,\mathrm{d}u\mathrm{d}v$

    所以

    \begin{align*}
        \text{原式}
         & =\frac{1}{2}\int_{-1}^1\,\mathrm{d}v\int_{-1}^1\frac{uv}{\sqrt{u+3}}\,\mathrm{d}u                   \\
         & =\frac{1}{2}\int_{-1}^1\,\mathrm{d}v\int_{-1}^1v(\sqrt{u+3}-\frac{3}{\sqrt{u+3}})\,\mathrm{d}u      \\
         & =\frac{1}{2}\int_{-1}^1v\left.\frac{2(u+3)^{\frac{3}{2}}}{3}-6\sqrt{u+3}\right|_{-1}^1\,\mathrm{d}v \\
         & =\frac{1}{2}(\frac{16}{3}-12-\frac{4\sqrt{2}}{3}+6\sqrt{2})\cdot 0                                  \\
         & =0
    \end{align*}
\end{sol}
(9) $\displaystyle\iint_{D} |xy|\,\mathrm{d}x\mathrm{d}y$, $D = \left\{(x,y):\,x^2 + y^2 \leq a^2\right\}$;
\begin{sol}
    令$\begin{cases}
            x=r\cos\theta \\
            y=r\sin\theta
        \end{cases}$, 所以$\mathrm{d}x\mathrm{d}y=r\,\mathrm{d}r\mathrm{d}\theta$

    原式为第一象限扇形面积的4倍，所以

    \begin{align*}
        \text{原式}
         & =4\int_0^{\frac{\pi}{2}}\sin\theta\cos\theta\,\mathrm{d}\theta\cdot\int_0^ar^2\,\mathrm{d}r \\
         & =4\cdot\left.\frac{\sin^2\theta}{2}\right|_0^{\frac{\pi}{2}}\cdot\frac{a^4}{4}              \\
         & =\frac{a^4}{2}
    \end{align*}
\end{sol}


3. 求下列曲线所围成的平面区域的面积:

(1) $x^2 + 2y^2 = 3$ 和 $xy = 1$ (不含原点部分);
\begin{sol}
    联立求解：$x^2+\frac{2}{x^2}=3$，解得第一象限内交点$(1,1)$, $(\sqrt{2},\frac{\sqrt{2}}{2})$

    因为围成不含原点的平面区域为关于原点中心对称的两部分，只需计算第一象限内的部分面积并乘2.

    在第一象限内：$\begin{cases}
            y=\sqrt{\frac{3-x^2}{2}} \\
            y=\frac{1}{x}
        \end{cases}$

    所以

    \begin{align*}
        \text{第一象限面积}
         & =\int_1^{\sqrt{2}}\sqrt{\frac{3-x^2}{2}}-\frac{1}{x}\,\mathrm{d}x                                             \\
         & =\left.\frac{\sqrt{2}x}{4}\sqrt{3-x^2}+\frac{3\sqrt{2}}{4}\arcsin\frac{x}{\sqrt{3}}-\ln x\right|_1^{\sqrt{2}}
    \end{align*}

    所以

    \begin{align*}
        \text{面积}
         & =\left.\frac{\sqrt{2}x}{2}\sqrt{3-x^2}+\frac{3\sqrt{2}}{2}\arcsin\frac{x}{\sqrt{3}}-\ln x^2\right|_1^{\sqrt{2}} \\
         & =(1+\frac{3\sqrt{2}}{2}\arcsin\sqrt{\frac{2}{3}}-\ln 2)-(1+\frac{3\sqrt{2}}{2}\arcsin\sqrt{\frac{1}{3}})        \\
         & =\frac{3\sqrt{2}}{2}(\arcsin\sqrt{\frac{2}{3}}-\arcsin\sqrt{\frac{1}{3}})-\ln 2
    \end{align*}
\end{sol}
(2) $(x - y)^2 + x^2 = a^2$ ($a > 0$);
\begin{sol}
    令$\begin{cases}
            u=x-y \\
            v=x
        \end{cases}$，所以$\frac{\partial(u,v)}{\partial(x,y)}=\begin{vmatrix}
            1 & -1 \\
            1 & 0
        \end{vmatrix}=1$, 所以$\mathrm{d}x\mathrm{d}y=\mathrm{d}u\mathrm{d}v$

    令$\begin{cases}
            u=r\cos\theta \\
            v=r\sin\theta
        \end{cases}$, 所以$\mathrm{d}u\mathrm{d}v=r\,\mathrm{d}r\mathrm{d}\theta$

    所以即求 $\int_0^{2\pi}\,\mathrm{d}\theta\int_0^a r\,\mathrm{d}r=\pi a^2$
\end{sol}
(3) 由直线 $x + y = a$, $x + y = b$, $y = kx$ 和 $y = mx$ ($0 < a < b$, $0 < k < m$) 围成的平面区域.
\begin{sol}
    令$\begin{cases}
            u=x+y \\
            v=\frac{y}{x}
        \end{cases}$, 所以$\begin{cases}
            a\leq u\leq b \\
            k\leq v\leq m
        \end{cases}$, $\frac{\partial(u,v)}{\partial(x,y)}=\begin{vmatrix}
            1              & 1           \\
            -\frac{y}{x^2} & \frac{1}{x}
        \end{vmatrix}=\frac{x+y}{x^2}=\frac{(1+v)^2}{u}$

    所以$\mathrm{d}x\mathrm{d}y=\frac{u}{(1+v)^2}\,\mathrm{d}u\mathrm{d}v$

    所以面积即

    \begin{align*}
        \int_k^m\,\mathrm{d}v\int_a^b\frac{u}{(1+v)^2}\,\mathrm{d}u
         & =\frac{b^2-a^2}{2}\int_k^m\frac{1}{(1+v)^2}\,\mathrm{d}v \\
         & =\frac{b^2-a^2}{2}(\frac{1}{1+k}-\frac{1}{1+m})          \\
         & =\frac{(a+b)(a-b)(k-m)}{2(1+k)(1+m)}
    \end{align*}
\end{sol}


4. 证明 $\displaystyle\iint_{x^2+y^2 \leq 1} \mathrm{e}^{x^2+y^2}\,\mathrm{d}x\mathrm{d}y \leq \left[ \int_{-\frac{\sqrt{\pi}}{2}}^{\frac{\sqrt{\pi}}{2}} \mathrm{e}^{x^2}\,\mathrm{d}x \right]^2$.
\begin{pf}
    LHS: 在单位圆内积分，记圆为$C$

    RHS: 在以原点为中心，边长为$\sqrt{\pi}$的正方形内积分，记正方形为$S$

    RHS$=\iint_{\substack{-\frac{\sqrt{\pi}}{2}\leq x\leq\frac{\sqrt{\pi}}{2}\\-\frac{\sqrt{\pi}}{2}\leq y\leq\frac{\sqrt{\pi}}{2}}}\mathrm{e}^{x^2+y^2}\,\mathrm{d}x\mathrm{d}y$

    所以$|C|=|S|=\pi$

    记 $C\cap S=D$, 所以$\begin{cases}
            C=C\setminus D+D=C\setminus S+D \\
            S=S\setminus D+D=S\setminus C+D
        \end{cases}$

    所以LHS$=\iint_{C\setminus S}\mathrm{e}^{x^2+y^2}\,\mathrm{d}x\mathrm{d}y+\iint_D\mathrm{e}^{x^2+y^2}\,\mathrm{d}x\mathrm{d}y$

    RHS$=\iint_{S\setminus C}\mathrm{e}^{x^2+y^2}\,\mathrm{d}x\mathrm{d}y+\iint_D\mathrm{e}^{x^2+y^2}\,\mathrm{d}x\mathrm{d}y$

    所以即证$\iint_{C\setminus S}\mathrm{e}^{x^2+y^2}\,\mathrm{d}x\mathrm{d}y\leq\iint_{S\setminus C}\mathrm{e}^{x^2+y^2}\,\mathrm{d}x\mathrm{d}y$

    在$C\setminus S$内，$x^2+y^2\leq 1$，所以$\iint_{C\setminus S}\mathrm{e}^{x^2+y^2}\,\mathrm{d}x\mathrm{d}y\leq\mathrm{e}|C\setminus S|$

    在$S\setminus C$内，$x^2+y^2\geq 1$，所以$\iint_{S\setminus C}\mathrm{e}^{x^2+y^2}\,\mathrm{d}x\mathrm{d}y\geq\mathrm{e}|S\setminus C|$

    又因为$|C\setminus S|=|S\setminus C|$, 所以

    \begin{align*}
        \iint_{C\setminus S}\mathrm{e}^{x^2+y^2}\,\mathrm{d}x\mathrm{d}y
         & \leq\mathrm{e}|C\setminus S|                                         \\
         & =\mathrm{e}|S\setminus C|                                            \\
         & \leq\iint_{S\setminus C}\mathrm{e}^{x^2+y^2}\,\mathrm{d}x\mathrm{d}y
    \end{align*}

    所以LHS$\leq$RHS

    Q.E.D.
\end{pf}


5. 设 $f(x)$ 在 $[0,1]$ 上连续, 证明 $\displaystyle\int_{0}^{1} \mathrm{e}^{f(x)}\,\mathrm{d}x \cdot \int_{0}^{1} \mathrm{e}^{-f(y)}\,\mathrm{d}y \geq 1$.
\begin{pf}
    令$\begin{cases}
            u(x)=\mathrm{e}^{\frac{f(x)}{2}} \\
            v(y)=\mathrm{e}^{-\frac{f(y)}{2}}
        \end{cases}$

    所以

    \begin{align*}
        \int_{0}^{1} \mathrm{e}^{f(x)}\,\mathrm{d}x \cdot \int_{0}^{1} \mathrm{e}^{-f(y)}\,\mathrm{d}y
         & =(\int_0^1u^2(x)\,\mathrm{d}x)\cdot(\int_0^1v^2(y)\,\mathrm{d}y)     \\
         & =(\int_0^1u^2(x)\,\mathrm{d}x)\cdot(\int_0^1v^2(x)\,\mathrm{d}x)     \\
         & \geq(\int_0^1u(x)v(x)\,\mathrm{d}x)^2                                \\
         & =(\int_0^1\mathrm{e}^{\frac{f(x)}{2}-\frac{f(x)}{2}}\,\mathrm{d}x)^2 \\
         & =(\int_0^1\mathrm{e}^0\,\mathrm{d}x)^2                               \\
         & =(\int_0^11\,\mathrm{d}x)^2                                          \\
         & =1
    \end{align*}

    Q.E.D.
\end{pf}


6. 设 $f(x)$ 为连续的奇函数, 证明 $\displaystyle\iint_{|x|+|y| \leq 1} \mathrm{e}^{f(x+y)}\,\mathrm{d}x\mathrm{d}y \geq 2$.
\begin{pf}
    令$\begin{cases}
            u=x+y \\
            v=x-y
        \end{cases}$, 所以$\begin{cases}
            -1\leq u\leq 1 \\
            -1\leq v\leq 1
        \end{cases}$, $\mathrm{d}x\mathrm{d}y=\frac{1}{2}\,\mathrm{d}u\mathrm{d}v$

    所以

    \begin{align*}
        \text{LHS}
         & =\frac{1}{2}\int_{-1}^1\,\mathrm{d}u\int_{-1}^1\mathrm{e}^{f(u)}\,\mathrm{d}v \\
         & =\int_{-1}^1\mathrm{e}^{f(u)}\,\mathrm{d}u                                    \\
    \end{align*}

    因为$f(x)$是连续的奇函数，所以

    \begin{align*}
        \text{LHS}
         & =\int_{-1}^1\mathrm{e}^{f(u)}\,\mathrm{d}u                                        \\
         & =\int_{-1}^0\mathrm{e}^{f(u)}\,\mathrm{d}u+\int_0^1\mathrm{e}^{f(u)}\,\mathrm{d}u \\
         & =\int_0^1\mathrm{e}^{-f(u)}\,\mathrm{d}u+\int_0^1\mathrm{e}^{f(u)}\,\mathrm{d}u   \\
         & =\int_0^1\mathrm{e}^{-f(u)}+\mathrm{e}^{f(u)}\,\mathrm{d}u                        \\
         & \geq\int_0^1 2\sqrt{\mathrm{e}^{-f(u)}\cdot\mathrm{e}^{f(u)}}\,\mathrm{d}u        \\
         & =\int_0^1 2\sqrt{\mathrm{e}^0}\,\mathrm{d}u                                       \\
         & =\int_0^1 2\,\mathrm{d}u                                                          \\
         & =2
    \end{align*}

    Q.E.D.
\end{pf}


7. 设 $f(t)$ 为连续函数, 求证
$$\iint_{D} f(x - y)\,\mathrm{d}x\mathrm{d}y = \int_{-A}^{A} f(t)\left(A - |t|\right)\,\mathrm{d}t$$
其中 $D$ 为: $|x| \leq \frac{A}{2}$, $|y| \leq \frac{A}{2}$, $A > 0$ 为常数.
\begin{pf}
    令$\begin{cases}
            u=x+y \\
            v=x-y
        \end{cases}$, 所以$|u|+|v|\leq A$, $\mathrm{d}x\mathrm{d}y=\frac{1}{2}\,\mathrm{d}u\mathrm{d}v$

    所以令$-A\leq v\leq A$, $-(A-|v|)\leq u\leq A-|v|$

    所以

    \begin{align*}
        \text{LHS}
         & =\frac{1}{2}\int_{-A}^A\,\mathrm{d}v\int_{-(A-|v|)}^{A-|v|}f(v)\,\mathrm{d}u \\
         & =\frac{1}{2}\int_{-A}^A2f(v)(A-|v|)\,\mathrm{d}v                             \\
         & =\int_{-A}^Af(v)(A-|v|)\,\mathrm{d}v
    \end{align*}

    将$v$替换成$t$，即得证

    Q.E.D.
\end{pf}

\end{document}